\color{unidarkgreen}{\subsection{Защита от неполнофазного режима }}
\color{black}

\begin{enumerate}[label=\arabic{section}.\arabic{subsection}.\arabic*,  align=left, labelsep=4pt, leftmargin=0pt, itemindent=57pt]

\item
Для контроля отключения всех фаз выключателя и ликвидации неполнофазного режима предусмотрена
функция защиты от неполнофазного режима (ЗНР).

\begin{figure}[H]
\centering
\includegraphics[width=0.88\textwidth,height=0.88\textheight,keepaspectratio]{\fbpath/021.0221 ЗНР/img/ZNR.pdf}
\captionof{figure}{Функциональная схема ЗНР ВН}\label{pict:ZNR}
\end{figure}

\item
Защита подключается к ТТ, установленным со стороны обмоток ВН силового трансформатора.

\item
Пуск ЗНР происходит при выполнении следующих условий:
\begin{itemize}
\item наличие сигнала срабатывания \textit{«ЗНФ В1 ВН в ЗНР»} или \textit{«ЗНФ В2 ВН в ЗНР»};
\item срабатывание измерительного органа, действующего по расчетному току нулевой последовательности или по отношению токов обратной и прямой последовательностей -- выбор осуществляется программным переключателем \textbf{«ИО ЗНР ВН»};
\item отключенное положение смежного выключателя (для схем с двумя выключателями).
\end{itemize}

\item
Наличие второго выключателя определяется программным переключателем \textbf{«Кол-во В ЗНР ВН»}. Для
полуторной схемы также реализован контроль отключенного положения выключателя В3 конфигурируемый
вход \textit{«В3 отключен в ЗНР»}.

\item
Порог срабатывания измерительного органа по расчетному току нулевой последовательности задается
уставкой \textbf{«3I0ср ЗНР ВН»}, по отношению токов обратной и прямой последовательностей -- \textbf{«I2/I1 ЗНР ВН»}.
Ввод в работу защиты от неполнофазного режима выключателей высокой стороны трансформатора
осуществляется с помощью программного переключателя \textbf{«ЗНР ВН»}, а выдержка времени срабатывания
задается уставкой \textbf{«Тср ЗНР ВН»}. 

\item
Оперативный ввод/вывод ЗНР ВН производится с помощью БУ (блоков управления функциями) \textit{«Вывод
ЗНР ВН»}. 


\small
\setlength{\extrarowheight}{0.05cm} % высота ячеек в таблице
\begin{longtable}
{|>{\centering\arraybackslash}m{4.5cm}|>{\centering\arraybackslash}m{4.5cm}|>{\centering\arraybackslash}m{3cm}|>{\centering\arraybackslash}m{3.5cm}|}
\caption{Перечень уставочных параметров ЗНР\hfill\vspace{-0.5\baselineskip}}\label{znr:tbl1} \\ 
\hhline{|-|-|-|-|}
\rowcolor{gray!30}
Наименование уставки & Диапазон & Шаг &  Значение по умолчанию \\ 
\hhline{|-|-|-|-|}
\endfirsthead

\caption*{\hspace{3pt}\emph{Продолжение таблицы \ref{znr:tbl1}\hfill\vspace{-0.5\baselineskip}}} \\ % сделано по ГОСТ 2.105 п.6.8.7
\hhline{|-|-|-|-|}
\rowcolor{gray!30}
Наименование уставки & Диапазон & Шаг &  Значение по умолчанию \\  
\hhline{|-|-|-|-|}
\endhead
\hhline{|-|-|-|-|}

\endfoot
\endlastfoot

%%%%%%%%%%%%содержимое%%%%%%%%%%%%%%

\end{longtable}
\setlength{\extrarowheight}{0.0cm} % возврат параметра высота ячеек в таблице
\normalsize










































\end{enumerate} 